\section*{الفصل \ref{ch:b3:innate-immunity} --- المناعة الفطرية والالتهاب}

\begin{solution}{exo:b3:innate-immunity:1}
الحواجز: كيراتين الجلد (ميكانيكي)، ومخاط الطرق الهوائية
وأهدابها (احتجاز وتصفية)، وحمض المعدة (يقتل الميكروبات
المبتلعة)، والليزوزيم والدفاعينات في الإفرازات (تحلّ
وتثقب)، \omterm{def:b3:microbiomes:symbiosis}{والميكروبيوم} المعايش (يشغل الكُوى). والخلايا:
العدلات (تقتل البكتيريا \omterm{def:b3:innate-immunity:killing}{بالبلعمة})، والبلعميات (تأكل وتنظف
وتطلق النذير)، والخلايا التغصنية (تحمل المستضد إلى العقد
اللمفية وتبدأ الاستجابات التكيّفية)، والخلايا البدينة (تطلق
الهيستامين وتُلهب)، والخلايا القاتلة الطبيعية (تقتل الخلايا
المصابة والمحوَّلة).
\end{solution}

\begin{solution}{exo:b3:innate-immunity:2}
أن يكون مشتركًا بين صنف كامل من الميكروبات، وضروريًا لها
(فلا يمكن التخلي عنه هربًا)، وغائبًا عن المضيف. \omterm{def:b3:bacteriology:envelope}{فعديد
السكريد الشحمي} --- TLR4؛ والفلاجيلين --- TLR5؛ وRNA مزدوج
الطاق --- TLR3 (وRIG-I في العصارة الخلوية)؛ وDNA ذو \omterm{def:b3:chromatin-epigenetics:methylation}{CpG} غير
الممثيل --- TLR9؛ وشظايا \omterm{def:b3:bacteriology:envelope}{الببتيدوغليكان} --- بروتينات NOD؛
وغلوكان $\beta$ --- الدكتين-1.
\end{solution}

\begin{solution}{exo:b3:innate-immunity:3}
الاحمرار: توسع الشرينات بالهيستامين وأكسيد النتريك
والبروستاغلاندينات. والحرارة: زيادة الجريان الدموي نفسها.
والتورم: رشح البلازما من الوُرَيدات بعد تقلص البطانة
(الهيستامين، وC3a، وC5a). والألم: البراديكينين
والبروستاغلاندين E$_{2}$ يحسّسان مستقبلات الألم، مع ضغط
التورم.
\end{solution}

\begin{solution}{exo:b3:innate-immunity:4}
التقليدي: ارتباط C1q بضد (أو \omterm{prop:b3:innate-immunity:systemic}{بالبروتين المتفاعل C}) على سطح.
واللكتيني: اللكتين الرابط للمانوز على سكريات ميكروبية.
والبديل: دوران C3 الذاتي العفوي وإيداع C3b على أي سطح غير
محمي. والنواتج: C3b يهيئ \omterm{def:b3:innate-immunity:killing}{للبلعمة} \omterm{def:b3:innate-immunity:complement}{بالأوبسنة}؛ وC3a وC5a
يُلهبان ويجنّدان؛ وC5b يبدأ \omterm{def:b3:innate-immunity:complement}{مركّب مهاجمة الغشاء} الذي يحلّ.
\end{solution}

\begin{solution}{exo:b3:innate-immunity:5}
$B = [k_{0}/(a-d)](e^{(a-d)t} - 1)$ عند $a - d = 0.1$: فعند \qty{60}{s}
يكون $20(e^{6} - 1) \approx 8000$؛ وعند \qty{90}{s} يكون $20(e^{9} - 1) \approx
1.6\times 10^{5}$. وعلى خلية مضيف: $k_{0}/(d - a) = 2/0.1 = 20$.
\end{solution}

\begin{solution}{exo:b3:innate-immunity:6}
تدحرج على السيليكتينات (لكتينات ترتبط بالسكريات)؛ وإشارة
كيموكينية تفعّل الإنتغرينات؛ والتصاق وثيق للإنتغرينات
بجزيئات الالتصاق البطانية (ICAM)؛ وعبور بين خلايا البطانة؛
وانجذاب كيميائي على امتداد التدرج. ومن دون إنتغرينات
$\beta_{2}$ تتدحرج العدلات ولا تتوقف ولا تعبر أبدًا:
فتتكدس في الدم (تعدادات عالية جدًا)، وتتكرر الأخماج بلا
تكوّن قيح، وتلتئم الجروح التئامًا رديئًا، ويتأخر انفصال
الحبل السري.
\end{solution}

\begin{solution}{exo:b3:innate-immunity:7}
الداء الحبيبي المزمن: \omterm{def:b3:innate-immunity:killing}{أكسيداز NADPH} معيب، فلا يُصنع في
الجسيم البلعمي فوق أكسيد ولا فوق أكسيد هدروجين ولا
هيبوكلوريت. والأخماج بكائنات تهدم فوق أكسيد الهدروجين الخاص
بها (إيجابية الكاتالاز): \emph{Staphylococcus aureus}،
و\emph{Aspergillus}، و\emph{Burkholderia}، و\emph{Serratia}،
و\emph{Nocardia}. وتتكون الأورام الحبيبية لأن البلعميات
تبتلع ميكروبات لا تقدر على قتلها فتسوّرها، مع تجنيد مستمر،
في آفة مزمنة.
\end{solution}

\begin{solution}{exo:b3:innate-immunity:8}
تعدّ المستقبلات المثبِّطة جزيئات MHC~I؛ وتعدّ المستقبلات
المفعِّلة روابط الإجهاد؛ وتموت الخلية إذا رجح التفعيل على
التثبيط. (أ) ففقدان MHC~I يرفع التثبيط: فتقتلها \omterm{def:b3:innate-immunity:innate}{الخلية
القاتلة الطبيعية} --- وذلك ثمن الهرب من اللمفاويات التائية.
(ب) ومع الاحتفاظ بجزيء MHC~I وعرض روابط الإجهاد، قد يغلب
التفعيل فتُقتل الخلية؛ ولذلك كثيرًا ما تنبذ الأورام روابط
الإجهاد أو تكبتها.
\end{solution}

\begin{solution}{exo:b3:innate-immunity:9}
البروتين النقي لا يحمل نمطًا ممرضًا، فالخلايا التغصنية التي
تلتقطه لا تنضج، وتعرضه بلا تحفيز مشارك، فتُطفأ اللمفاويات
التائية التي تتعرفه أو تتجاهله. \omterm{prop:b3:innate-immunity:bridge}{والمادة المساعدة} --- شبّ، أو
شحم، أو رابطة \omterm{def:b3:innate-immunity:prr}{لمستقبل شبيه بمستقبل Toll} --- تثير المستقبلات
النمطية، فتُنضج \omterm{def:b3:innate-immunity:innate}{الخلية التغصنية} وتوفر الإشارة الثانية. وتلك
حجة جانواي: فالجهاز التكيّفي لا يستجيب للمستضد إلا حين يقول
الجهاز الفطري إن المستضد جاء مع ميكروب.
\end{solution}

\begin{solution}{exo:b3:innate-immunity:10}
بلا حمّى، $48$ تضاعفًا: $10^{4}\times 2^{48} \approx 3\times
10^{18}$ (وهو عبث يبيّن أن القتل، لا النمو، هو الذي يقرر
النتيجة). ومع الحمّى، $32$ تضاعفًا: $4\times 10^{13}$ --- أي بعامل
$2^{16} \approx 65\,000$ من البكتيريا أقل تلقاها العدلات.
والكلفة: $3\times 0.12\times 8 = \qty{2.9}{MJ}$ في اليوم، أي
وجبة كبيرة. وهي تستحق حين يُبطأ الممرض ويكون للمضيف
احتياطي؛ ولا تستحق حين ينمو الممرض نموًا جيدًا عند
\qty{40}{\celsius}، أو حين يكون المضيف جائعًا، أو حين تقارب
الحمّى نفسها حرارات خطرة.
\end{solution}

\begin{solution}{exo:b3:innate-immunity:11}
ربط العامل H يجنّد منظِّم المضيف نفسه: فيرتفع $d$ فوق $a$
وتخمد العروة على الميكروب كما تخمد على خلية مضيف. وشطر C3b
يزيل المحوّلة المودعة: وهو أيضًا ارتفاع في $d$. والمحفظة
تقدم سطحًا تتكون عليه المحوّلة تكوّنًا رديئًا، فتخفض $a$،
وتخفي C3b المودع عن مستقبلات البالعات. وأرخصها: بروتين سطحي
واحد يربط العامل H، الذي يوفره المضيف مجانًا ويقوم بالعمل.
\end{solution}

\begin{solution}{exo:b3:innate-immunity:12}
عند الشظية يشمل التوسع والرشح والتجنيد بضعة ميليلترات من
الأوعية ويوصل الدفاع حيث يلزم؛ وأما الوسائط نفسها إذا
أُطلقت في الدوران كله فتوسّع كل وعاء وترشّح كل شعيرة،
فينهار الضغط ويُطلق التخثر في كل مكان --- فالفيزيولوجيا
واحدة، والمقياس قاتل. وTNF يقود \omterm{def:b3:innate-immunity:inflammation}{الالتهاب} الزليلي في \omterm{def:b3:innate-immunity:inflammation}{التهاب}
المفاصل الرثوي، فحجبه يخفف المرض؛ غير أن TNF هو أيضًا ما
يبقي البلعميات منظَّمة في الأورام الحبيبية التي تحتوي عصيات
السل الهاجعة، وحجبه يطلق السل الكامن.
\end{solution}

\begin{solution}{pb:b3:innate-immunity:1}
\textbf{1.} الساعة الواحدة $1.5$ تضاعف: $10^{4}\times 2^{1.5} \approx
2.8\times 10^{4}$.
\textbf{2.} تقتل $2\times 10^{4}$ \omterm{def:b3:innate-immunity:innate}{عدلة} ما يصل إلى $4\times 10^{5}$؛
ولا تنمو البكتيريا إلا من $2.8\times 10^{4}$ إلى $8\times 10^{4}$:
أي إن القتل يفوق النمو خمسة أضعاف.
\textbf{3.} الساعة 2: $8\times 10^{4}$ قبل القتل، ولا شيء بعده.
فتبلغ الجمهرة ذروتها عند نحو $8\times 10^{4}$ قرب \qty{2}{h}
وتكون صفرًا من ثَمّ؛ ولا يبقى منها شيء عند \qty{6}{h}.
\textbf{4.} الوصول عند \qty{3}{h}: $10^{4}\times 2^{4.5} = 2.3\times
10^{5}$ عندئذ. والساعة 4: $6.4\times 10^{5} - 4\times 10^{5} = 2.4\times
10^{5}$؛ والساعة 5: $6.8\times 10^{5} - 4\times 10^{5} = 2.8\times 10^{5}$؛
والساعة 6: $7.9\times 10^{5} - 4\times 10^{5} = 3.9\times 10^{5}$ وفي
ارتفاع --- فلم تعد العدلات تجاري، ويتكون خراج. فساعتا تأخر
تحيلان الشفاء خمجًا مزمنًا.
\textbf{5.} $8\times 10^{4}/20 = 4000$ \omterm{def:b3:innate-immunity:innate}{عدلة} تموت وهي تقتل، مع
$2\times 10^{4}$ وصلت وتموت على كل حال في غضون يوم. والقيح عدلات
ميتة وبكتيريا ميتة ونسيج متميّع؛ وتصفّيه البلعميات بأكل
الجثث، أو لا بد من تصريفه.
\textbf{6.} الوافدات $2000$ في الساعة، تقتل $4\times 10^{4}$: فالساعة 2،
$8\times 10^{4} - 4\times 10^{4} = 4\times 10^{4}$؛ والساعة 3، $1.1\times
10^{5} - 4\times 10^{4} = 7\times 10^{4}$؛ والساعة 4، $2\times 10^{5} -
4\times 10^{4} = 1.6\times 10^{5}$ --- في نمو بلا حد. فالمريض لا
يقدر على احتواء ملقَّحة تافهة، فلا بد أن تتولى الصادات
القتل من أول علامة.
\textbf{7.} $B = 12.5(e^{0.08t} - 1)$: فعند \qty{30}{s}، $125$؛ وعند \qty{60}{s}،
$1500$؛ وعند \qty{120}{s}، $1.8\times 10^{5}$.
\textbf{8.} $e^{0.08t} = 1.6\times 10^{5}$: أي $t = 12/0.08 = \qty{150}{s}$.
\textbf{9.} $k_{0}/(d - a) = 10$.
\textbf{10.} $B(120) = 33(e^{3.6} - 1) \approx 1200$، أي أقل $150$ مرة.
فالمحفظة تشتري دقائق --- وقتًا للتكاثر، وحمايةً إلى أن تصل
الأضداد فتبدأ السبيل التقليدي على المحفظة نفسها.
\textbf{11.} $10^{4}\times 2\times 10^{6} = 2\times 10^{10}$ من C3b: أي
جزء من مليون من C3 في البلازما.
\textbf{12.} من دون C3 تتوقف السُبل الثلاثة كلها عند خطوتها
المشتركة: فلا \omterm{def:b3:innate-immunity:complement}{أوبسنة}، ولا ذيفانات تأقية، ولا حلّ --- أي
أخماج شديدة متكررة ببكتيريا محفظة. ومن دون C9 لا يُفقد إلا
المسام الطرفي؛ \omterm{def:b3:innate-immunity:complement}{فالأوبسنة} \omterm{def:b3:innate-immunity:inflammation}{والالتهاب} يعملان، ويكون النمط
الظاهري أخماج \emph{Neisseria} المتكررة، وهو الجنس الوحيد
الذي تضبطه البالعات ضبطًا رديئًا ويضبطه الحلّ ضبطًا جيدًا.
\textbf{13.} $2\times 0.24\times 8 = \qty{3.8}{MJ}$، أي نحو \qty{100}{g}
من الشحم.
\textbf{14.} تضاعف واحد في الساعة: $2\times 10^{4}$ بدل
$2.8\times 10^{4}$، أي بعامل $1.4$.
\textbf{15.} الساعة 2: $4\times 10^{4}$ قبل القتل، ولا شيء بعده ---
فتزول عند \qty{2}{h}، كما في السابق؛ وإسهام الحمّى يهم حين
يكون مدد العدلات حديًّا، كما في السؤالين 4 و6.
\textbf{16.} $199\,\text{mg/L}\times 3\,\text{L} = \qty{0.6}{g}$؛ و$0.6/
\num{115000} = 5.2\times 10^{-6}$ مول $= 3.1\times 10^{18}$ جزيء؛
وعلى مدى \qty{172800}{s}، أي $1.8\times 10^{13}$ في الثانية.
\textbf{17.} تخفض نقطة الضبط راجعةً نحو \qty{37}{\celsius}، فيشعر
المريض بتحسن، وتتضاعف البكتيريا أسرع قليلًا؛ وفي هذا
النموذج تظل العدلات غالبة، والشواهد على أن خافضات الحرارة
تسوئ الأخماج العادية ضعيفة.
\textbf{18.} فوق نحو \qty{41}{\celsius} تبدأ البروتينات بالتمسخ،
وتخفق الإنزيمات، وتشتط العصبونات (اختلاجات)؛ و\qty{42}{\celsius}
قاتلة. وتُبقى نقطة الضبط دون ذلك بمبرّدات داخلية ---
الإنترلوكين-10، والفازوبريسين، والقشرانيات السكرية --- وبسقف
فعل البروستاغلاندين في الوطاء.
\textbf{19.} $10^{9}\times 3\times 10^{6} = 3\times 10^{15}$ جزيء
$= 5\times 10^{-9}$ مول في \qty{5}{L}: أي \qty{1e-9}{mol/L}.
\textbf{20.} مساوٍ للتركيز المشبِع: أي إن كل \omterm{def:b3:innate-immunity:innate}{بلعمية} في الجسم
مفعَّلة تفعيلًا أقصى دفعة واحدة.
\textbf{21.} $10^{11}\times 1000\times 3600 = 3.6\times 10^{17}$ جزيء
$= 6\times 10^{-7}$ مول؛ وفي \qty{15}{L}، \qty{4e-8}{mol/L} --- أي أربعة
آلاف ضعف التركيز الفعال.
\textbf{22.} يتوسع كل شرين ويرشح كل وُرَيد: فيتجمع حجم الدم في
الأوعية المتسعة ويهرب إلى النسج، فيهبط الضغط؛ وتكف
الكليتان، إذ ينقص إرواؤهما، عن الرشح؛ ويسد التخثر المطلَق
في آلاف الشعيرات تلك الشعيرات ويستهلك عوامل التخثر.
\textbf{23.} سلبيات الغرام المقتولة تطلق \omterm{def:b3:bacteriology:envelope}{عديد السكريد الشحمي}
كله دفعة واحدة، وهو جرعة لمستقبل TLR4، فقد تسوء طفرة
السيتوكينات ساعات. ويخفق حجب TNF لأن الشلال يكون، وقت
الصدمة، قد تجاوز TNF --- إذ يجري الإنترلوكين-1
والإنترلوكين-6 والتخثر فعلًا --- ولأن TNF لازم لضبط
البكتيريا: فالمشكلة في التوقيت لا في الهدف.
\textbf{24.} \omterm{def:b3:bacteriology:envelope}{عديد السكريد الشحمي} المحقون لا يفعل شيئًا بالفأر
العديم TLR4، الذي ينجو من جرعات قاتلة للفئران السوية؛ وأما
خمج حي بسلبيات الغرام فيقتله، لأنه لا يقدر على كشف
البكتيريا باكرًا واستنفار \omterm{def:b3:innate-immunity:inflammation}{الالتهاب} الذي يحتويها. وتنحل
المفارقة: فالكاشف الذي يسبب الصدمة هو الكاشف الذي يوفر
الدفاع.
\textbf{25.} لا تبقى بكتيريا عند \qty{6}{h} (قُتلت كلها عند
\qty{2}{h})؛ ونحو $1.8\times 10^{5}$ من C3b لكل جرثوم عند
\qty{120}{s}؛ ويكلف يومان من الحمّى \qty{3.8}{MJ}، أي نحو
\qty{100}{g} من الشحم.
\end{solution}
